\subsection{分割格 prime-slice 极小刚性与 Schur 单变量主核}\label{subsec:conclusion-partition-primeslice-schur-singlekernel}
在集合分割格的 prime-slice meet 算术模型与 Schur 断层的 Frobenius/Jacobi--Trudi 描述都已建立之后，还可把二者各自推进为两条严格的唯一性原则。其一，极小素数预算不只给出 meet--\(\gcd\) 模型的存在性，而且把整个模型刚性化为“除边--素数重命名外别无自由度”，并使固定分割的全部 squarefree 编码形成一族精确因子化的有限温配分函数，其零温极限唯一集中到最小骨架编码。其二，Schur 断层在固定分辨率上的全部高维通道并非独立对象：单行通道 \(T_{q,(q)}(m)=h_q(D_m)\) 已足以恢复纤维多重度多重集，从而恢复全部 Schur 通道，而整个通道塔又恰由一条单变量核通过 Cauchy 外化完全生成。

\begin{theorem}[极小 meet--算术模型的唯一性与 prime-slice 刚性]\label{thm:conclusion-partition-primeslice-minimal-uniqueness}
设 \(n\ge 3\)。沿用定义 \ref{def:pom-partition-lattice}、定义 \ref{def:pom-squarefree-edge-code} 与命题 \ref{prop:pom-closure-fixedpoints-partitions} 的记号，记
\[
E_n:=\binom{[n]}2,
\qquad
F_\pi:=\bigcup_{B\in\pi}\binom{B}{2}
\qquad
\bigl(\pi\in\mathcal P([n])\bigr).
\]
设 \(S\) 为有限素数集，且
\[
\Phi:\mathcal P([n])\hookrightarrow \NN_{\mathrm{sf}}(S)
\]
为满足
\[
\Phi(\pi\wedge\sigma)=\gcd\bigl(\Phi(\pi),\Phi(\sigma)\bigr)
\qquad
\bigl(\forall\,\pi,\sigma\in\mathcal P([n])\bigr)
\]
的单射。若
\[
|S|=\binom{n}{2},
\]
则存在唯一双射
\[
\eta:E_n\overset{\sim}{\longrightarrow}S
\]
使得对任意 \(\pi\in\mathcal P([n])\) 都有
\[
\boxed{\;
\Phi(\pi)=\prod_{e\in F_\pi}\eta(e).\;}
\]
换言之，在极小素数预算下，满足 meet--\(\gcd\) 严格交换的 prime-slice 模型除边--素数重命名外没有任何额外自由度。
\end{theorem}
\begin{proof}
记 \(\hat 0\in\mathcal P([n])\) 为离散分割，并对每条边 \(e\in E_n\) 记 \(\alpha_e\) 为只合并 \(e\) 两端点的原子分割。由定理 \ref{thm:pom-meet-gcd-prime-budget} 的证明机制，若记
\[
d:=\Phi(\hat 0),
\]
则对每个 \(e\in E_n\) 都可写成
\[
\Phi(\alpha_e)=d\,a_e,
\]
其中 \(a_e\neq 1\)，并且 \(\{a_e\}_{e\in E_n}\) 两两互素。又由于 \(\Phi(\alpha_e)\in \NN_{\mathrm{sf}}(S)\) 为平方自由整数，必有
\[
\gcd(d,a_e)=1
\qquad
\bigl(\forall\,e\in E_n\bigr).
\]

现在利用极小预算 \(|S|=|E_n|\)。因为每个 \(a_e\neq 1\) 且彼此互素，所以每个 \(a_e\) 至少贡献一个新的素因子；但可用素数总数恰为 \(|E_n|\)，故每个 \(a_e\) 只能恰含一个素因子，从而必为素数，且
\[
\{a_e:\ e\in E_n\}=S.
\]
若 \(d\neq 1\)，取任意素数 \(p\mid d\)。由于 \(p\in S\)，存在唯一 \(e_0\in E_n\) 使 \(a_{e_0}=p\)。于是
\[
p^2\mid d\,a_{e_0}=\Phi(\alpha_{e_0}),
\]
这与 \(\Phi(\alpha_{e_0})\) 平方自由矛盾。故必有
\[
d=1.
\]
于是
\[
\Phi(\alpha_e)=a_e
\qquad
\bigl(e\in E_n\bigr)
\]
本身就是一组与边集 \(E_n\) 一一对应的互异素数。定义
\[
\eta(e):=\Phi(\alpha_e)=a_e.
\]

对任意 \(\pi\in\mathcal P([n])\) 与任意 \(e\in E_n\)，有
\[
\alpha_e\le \pi
\iff
\alpha_e\wedge\pi=\alpha_e
\iff
\gcd\bigl(\Phi(\alpha_e),\Phi(\pi)\bigr)=\Phi(\alpha_e)
\iff
\eta(e)\mid \Phi(\pi).
\]
另一方面，由 \(F_\pi\) 的定义，\(\alpha_e\le \pi\) 当且仅当 \(e\in F_\pi\)。因此 \(\Phi(\pi)\) 的素因子集合恰为
\[
\{\eta(e):\ e\in F_\pi\}.
\]
又由于 \(\Phi(\pi)\in \NN_{\mathrm{sf}}(S)\) 平方自由，便得到
\[
\Phi(\pi)=\prod_{e\in F_\pi}\eta(e).
\]
唯一性显然，因为对每个 \(e\in E_n\) 都有
\[
\eta(e)=\Phi(\alpha_e),
\]
故 \(\eta\) 已由原子像完全决定。证毕。
\end{proof}
\leanverified{paper\_conclusion\_partition\_primeslice\_minimal\_uniqueness}

\begin{corollary}[极小 prime-slice 编码同时刚性化偏序、原子支撑与块内配对计数]\label{cor:conclusion-partition-primeslice-divisibility-valuation-paircount}
在定理 \ref{thm:conclusion-partition-primeslice-minimal-uniqueness} 的设定下，对任意 \(\pi,\sigma\in\mathcal P([n])\) 都有
\[
\pi\le \sigma
\iff
\Phi(\pi)\mid \Phi(\sigma).
\]
并且对任意边 \(e\in E_n\)，有
\[
v_{\eta(e)}\bigl(\Phi(\pi)\bigr)=\mathbf 1_{\{e\in F_\pi\}},
\]
以及
\[
\Omega\bigl(\Phi(\pi)\bigr)
=
|F_\pi|
=
\sum_{B\in\pi}\binom{|B|}{2}.
\]
因此，在极小素数预算下，分割格的偏序、原子支撑与块内配对总数都被同一个平方自由整数 \(\Phi(\pi)\) 完全而唯一地承载。
\end{corollary}
\begin{proof}
由定理 \ref{thm:conclusion-partition-primeslice-minimal-uniqueness}，
\[
\Phi(\pi)=\prod_{e\in F_\pi}\eta(e),
\qquad
\Phi(\sigma)=\prod_{e\in F_\sigma}\eta(e).
\]
由于 \(\eta(e)\) 两两互异且 \(\Phi(\pi),\Phi(\sigma)\) 平方自由，故
\[
\Phi(\pi)\mid \Phi(\sigma)
\iff
F_\pi\subseteq F_\sigma.
\]
而正文中的 partition-closure 判据给出
\[
F_\pi\subseteq F_\sigma
\iff
\pi\le \sigma.
\]
这就得到第一式。

第二式同样由平方自由分解立即得出：素数 \(\eta(e)\) 在 \(\Phi(\pi)\) 中出现当且仅当 \(e\in F_\pi\)，且其指数恰为 \(1\)。最后，因为 \(\Phi(\pi)\) 平方自由，\(\Omega(\Phi(\pi))\) 正是其不同素因子个数，故
\[
\Omega\bigl(\Phi(\pi)\bigr)=|F_\pi|.
\]
而
\[
F_\pi=\bigsqcup_{B\in\pi}\binom{B}{2},
\]
故
\[
|F_\pi|=\sum_{B\in\pi}\binom{|B|}{2}.
\]
证毕。
\end{proof}


\begin{theorem}[固定分割 squarefree 编码的配分函数分解与零温唯一基态]\label{thm:conclusion-partition-squarefree-partition-function-groundstate}
沿用定义 \ref{def:pom-squarefree-edge-code} 与定理 \ref{thm:pom-partition-skeleton-compression} 的记号。固定分割
\[
\pi=\{B_1,\dots,B_k\}\in\mathcal P([n]),
\]
并定义
\[
\mathcal Z_\pi(s)
:=
\sum_{\substack{F\subseteq E_n\\ \mathrm{cl}(F)=F_\pi}}
\mathsf{Code}(F)^{-s}
\qquad
\bigl(\Re s\ge 0\bigr).
\]
对任意非空块 \(B\subseteq[n]\)，再定义局部连通图配分函数
\[
\mathcal C_B(s)
:=
\sum_{\substack{G\subseteq \binom{B}{2}\\ (B,G)\ \mathrm{连通}}}
\ \prod_{e\in G}p_e^{-s}.
\]
则有严格因子分解
\[
\boxed{\;
\mathcal Z_\pi(s)=\prod_{t=1}^{k}\mathcal C_{B_t}(s).\;}
\]
特别地，
\[
\boxed{\;
\mathcal Z_\pi(0)=\prod_{t=1}^{k}c_{|B_t|},\;}
\]
其中 \(c_r\) 为 \(r\) 个标号顶点上的连通图总数。

进一步，定义能隙
\[
\gamma_\pi
:=
\inf\left\{
\log\frac{\mathsf{Code}(F)}{\mathsf{Code}_{\min}(\pi)}:
\ \mathrm{cl}(F)=F_\pi,\ 
\mathsf{Code}(F)>\mathsf{Code}_{\min}(\pi)
\right\}\in(0,\infty],
\]
则当 \(\Re s\to+\infty\) 时有
\[
\boxed{\;
\mathcal Z_\pi(s)
=
\mathsf{Code}_{\min}(\pi)^{-s}
\Bigl(1+O\!\bigl(e^{-\gamma_\pi\,\Re s}\bigr)\Bigr).\;}
\]
当 \(\gamma_\pi=+\infty\) 时，上式中的误差项恒为零；这恰对应可行编码唯一的情形。
\end{theorem}
\begin{proof}
由命题 \ref{prop:pom-closure-fixedpoints-partitions}，条件
\[
\mathrm{cl}(F)=F_\pi
\]
等价于以下两条同时成立：
\begin{enumerate}
  \item \(F\) 不含跨不同块的边；
  \item 对每个 \(t=1,\dots,k\)，限制图 \(\bigl(B_t,F\cap\binom{B_t}{2}\bigr)\) 连通。
\end{enumerate}
因此任意可行边集 \(F\) 都唯一分解为
\[
F=G_1\sqcup\cdots\sqcup G_k,
\qquad
G_t\subseteq \binom{B_t}{2},
\qquad
(B_t,G_t)\ \text{连通}.
\]
再由编码的乘法性，
\[
\mathsf{Code}(F)^{-s}
=
\prod_{t=1}^{k}
\left(\prod_{e\in G_t}p_e^{-s}\right).
\]
对各块分别求和便得到
\[
\mathcal Z_\pi(s)=\prod_{t=1}^{k}\mathcal C_{B_t}(s).
\]

取 \(s=0\) 时，每一项权重均为 \(1\)，故 \(\mathcal C_{B_t}(0)\) 恰为顶点集 \(B_t\) 上连通图的总数；该数量只依赖于 \(|B_t|\)，记为 \(c_{|B_t|}\)，遂得
\[
\mathcal Z_\pi(0)=\prod_{t=1}^{k}c_{|B_t|}.
\]

最后，由定理 \ref{thm:pom-partition-skeleton-compression}，\(\mathsf{Code}_{\min}(\pi)\) 存在且唯一。设
\[
\mathcal F_\pi:=\{F\subseteq E_n:\ \mathrm{cl}(F)=F_\pi\}.
\]
则 \(\mathcal F_\pi\) 为有限集，并且在 \(\mathcal F_\pi\) 中存在唯一最小编码 \(F_{\min}\) 满足
\[
\mathsf{Code}(F_{\min})=\mathsf{Code}_{\min}(\pi).
\]
于是
\[
\mathcal Z_\pi(s)
=
\mathsf{Code}_{\min}(\pi)^{-s}
\left[
1+
\sum_{F\in\mathcal F_\pi\setminus\{F_{\min}\}}
\exp\!\left(
-s\log\frac{\mathsf{Code}(F)}{\mathsf{Code}_{\min}(\pi)}
\right)
\right].
\]
对任意 \(F\neq F_{\min}\)，括号中的对数均至少为 \(\gamma_\pi\)。由于 \(\mathcal F_\pi\) 有限，便有一致估计
\[
\sum_{F\in\mathcal F_\pi\setminus\{F_{\min}\}}
\left|
\exp\!\left(
-s\log\frac{\mathsf{Code}(F)}{\mathsf{Code}_{\min}(\pi)}
\right)
\right|
=
O\!\bigl(e^{-\gamma_\pi\,\Re s}\bigr).
\]
这就得到所示零温展开。若 \(\mathcal F_\pi=\{F_{\min}\}\)，则 \(\gamma_\pi=+\infty\)，此时 \(\mathcal Z_\pi(s)=\mathsf{Code}_{\min}(\pi)^{-s}\) 而误差项恒为零。证毕。
\end{proof}

\begin{corollary}[固定分割 squarefree 编码 Gibbs 律的零温单态塌缩]\label{cor:conclusion-partition-squarefree-zero-temp-gibbs-collapse}
在定理 \ref{thm:conclusion-partition-squarefree-partition-function-groundstate} 的设定下，定义可行边集上的 Gibbs 测度
\[
\mu_{\pi,s}(F)
:=
\frac{\mathsf{Code}(F)^{-s}}{\mathcal Z_\pi(s)}
\qquad
\bigl(F\subseteq E_n,\ \mathrm{cl}(F)=F_\pi\bigr).
\]
则零温自由能满足
\[
-\lim_{s\to+\infty}\frac1s\log \mathcal Z_\pi(s)
=
\log \mathsf{Code}_{\min}(\pi),
\]
并且
\[
\mu_{\pi,s}\Longrightarrow \delta_{F_{\min}(\pi)}
\qquad
(s\to+\infty).
\]
换言之，固定分割的 squarefree 编码热力学在零温极限下严格塌缩到唯一最小骨架。
\end{corollary}
\begin{proof}
由定理 \ref{thm:conclusion-partition-squarefree-partition-function-groundstate}，
\[
\mathcal Z_\pi(s)
=
\mathsf{Code}_{\min}(\pi)^{-s}
\Bigl(1+O\!\bigl(e^{-\gamma_\pi s}\bigr)\Bigr)
\qquad
(s\to+\infty),
\]
故对数后除以 \(s\) 即得自由能极限。

再对任意 \(F\neq F_{\min}(\pi)\)，有
\[
\mu_{\pi,s}(F)
=
\frac{\exp\!\left(-s\log\frac{\mathsf{Code}(F)}{\mathsf{Code}_{\min}(\pi)}\right)}
{1+\sum_{G\neq F_{\min}(\pi)}\exp\!\left(-s\log\frac{\mathsf{Code}(G)}{\mathsf{Code}_{\min}(\pi)}\right)}
\longrightarrow 0,
\]
因为分子含有严格正的能隙。于是
\[
\mu_{\pi,s}\bigl(F_{\min}(\pi)\bigr)\longrightarrow 1,
\]
从而 \(\mu_{\pi,s}\) 弱收敛到 \(\delta_{F_{\min}(\pi)}\)。证毕。
\end{proof}


\begin{theorem}[单行 Schur 通道对纤维多重度谱与全通道塔的完备恢复]\label{thm:conclusion-schur-single-row-complete-recovery}
固定分辨率 \(m\)。沿用推论 \ref{cor:pom-schur-trace-frobenius-character-identification} 与命题 \ref{prop:pom-schur-trace-jacobi-trudi-euler-primitive-fiber} 的记号，记
\[
N_m:=|X_m|,
\qquad
D_m:=\mathrm{diag}\bigl(d_m(x)\bigr)_{x\in X_m}.
\]
并补充约定
\[
T_{1,(1)}(m):=h_1(D_m)=\sum_{x\in X_m}d_m(x).
\]
则有限序列
\[
\bigl(
T_{1,(1)}(m),\,
T_{2,(2)}(m),\,
\dots,\,
T_{N_m,(N_m)}(m)
\bigr)
\]
已经唯一决定整个纤维多重度多重集
\[
\{d_m(x):x\in X_m\},
\]
从而唯一决定全部 Schur 通道迹
\[
T_{q,\lambda}(m)
\qquad
\bigl(q\ge 1,\ \lambda\vdash q\bigr).
\]
更精确地，若记 \(e_r(D_m)\) 为 \(D_m\) 对角元的第 \(r\) 个基本对称函数，并约定
\[
e_0(D_m)=1,
\qquad
h_0(D_m)=1,
\qquad
h_j(D_m)=0\ \ (j<0),
\]
则对任意 \(q\ge 1\) 都有递推
\[
\boxed{\;
\sum_{r=0}^{N_m}
(-1)^r e_r(D_m)\,h_{q-r}(D_m)=0.\;}
\]
其中
\[
h_q(D_m)=T_{q,(q)}(m).
\]
因此多重度特征多项式
\[
\chi_m(T):=\prod_{x\in X_m}(T-d_m(x))
=
T^{N_m}-e_1(D_m)T^{N_m-1}+\cdots+(-1)^{N_m}e_{N_m}(D_m)
\]
由前 \(N_m\) 个单行通道值唯一决定。
\end{theorem}
\begin{proof}
由推论 \ref{cor:pom-schur-trace-frobenius-character-identification}（对 \(q\ge 2\)）以及上述对 \(q=1\) 的约定，
\[
T_{q,(q)}(m)=h_q(D_m)
\qquad
\bigl(q\ge 1\bigr).
\]
再由命题 \ref{prop:pom-schur-trace-jacobi-trudi-euler-primitive-fiber}，
\[
H_m(z):=\sum_{q\ge 0}h_q(D_m)z^q
=
\prod_{x\in X_m}\frac{1}{1-d_m(x)z}.
\]
于是
\[
\left(
\sum_{r=0}^{N_m}(-1)^r e_r(D_m)z^r
\right)
\left(
\sum_{q\ge 0}h_q(D_m)z^q
\right)
=
\prod_{x\in X_m}(1-d_m(x)z)\prod_{x\in X_m}(1-d_m(x)z)^{-1}
=
1.
\]
比较 \(z^q\) 系数便得
\[
\sum_{r=0}^{N_m}(-1)^r e_r(D_m)\,h_{q-r}(D_m)=0
\qquad(q\ge 1),
\]
其中约定 \(h_j(D_m)=0\)（\(j<0\)）。

当 \(q=1,\dots,N_m\) 时，上式形成关于
\[
e_1(D_m),\dots,e_{N_m}(D_m)
\]
的三角递推系统。由于 \(h_0(D_m)=1\) 已知，而
\[
h_q(D_m)=T_{q,(q)}(m)
\qquad
\bigl(1\le q\le N_m\bigr)
\]
亦已知，故可依次唯一求出全部 \(e_r(D_m)\)。于是 \(\chi_m(T)\) 的系数唯一确定，而它的根多重集正是
\[
\{d_m(x):x\in X_m\}.
\]
这就恢复了整个纤维多重度多重集。

最后，再由推论 \ref{cor:pom-schur-trace-frobenius-character-identification}，
\[
T_{q,\lambda}(m)=f^\lambda s_\lambda(D_m)
\qquad
\bigl(\lambda\vdash q\bigr).
\]
既然 \(D_m\) 的对角元多重集已知，全部 \(s_\lambda(D_m)\) 乃至全部 \(T_{q,\lambda}(m)\) 便唯一确定。证毕。
\end{proof}

\begin{theorem}[Schur 断层主核的 Cauchy 外化与单变量生成]\label{thm:conclusion-schur-master-kernel-cauchy-externalization}
沿用定理 \ref{thm:conclusion-schur-single-row-complete-recovery} 的记号，并定义单变量主核
\[
H_m(z):=\sum_{q\ge 0}h_q(D_m)z^q
=
1+\sum_{q\ge 1}T_{q,(q)}(m)\,z^q.
\]
再对辅助变量组 \(\mathbf y=(y_1,y_2,\dots)\) 定义 Schur 主核
\[
\mathcal K_m(\mathbf y)
:=
1+\sum_{q\ge 1}\ \sum_{\lambda\vdash q}
\frac{T_{q,\lambda}(m)}{f^\lambda}\,s_\lambda(\mathbf y).
\]
则有严格恒等式
\[
\boxed{\;
\mathcal K_m(\mathbf y)
=
\prod_{x\in X_m}\prod_{j\ge 1}\frac{1}{1-d_m(x)y_j}
=
\prod_{j\ge 1}H_m(y_j).\;}
\]
因此：
\begin{enumerate}
  \item 整个 Schur 通道塔由单变量核 \(H_m\) 通过 Cauchy 外化 \(\prod_{j\ge 1}H_m(y_j)\) 完全生成；
  \item 对任意 \(\lambda\vdash q\)，存在仅依赖于 \(\lambda\) 的整系数多项式 \(\mathfrak S_\lambda\in\ZZ[U_1,\dots,U_q]\) 使
  \[
  \boxed{\;
  \frac{T_{q,\lambda}(m)}{f^\lambda}
  =
  \mathfrak S_\lambda\bigl(
  T_{1,(1)}(m),\dots,T_{q,(q)}(m)
  \bigr).\;}
  \]
  \item 把 \(\mathbf y\) 特化为 \((z,0,0,\dots)\) 时，恢复
  \[
  \boxed{\;
  \mathcal K_m(z,0,0,\dots)=H_m(z).\;}
  \]
\end{enumerate}
\end{theorem}
\begin{proof}
对 \(q\ge 2\)，由推论 \ref{cor:pom-schur-trace-frobenius-character-identification} 有
\[
\frac{T_{q,\lambda}(m)}{f^\lambda}=s_\lambda(D_m).
\]
而在唯一的 \(q=1,\lambda=(1)\) 情形，定理 \ref{thm:conclusion-schur-single-row-complete-recovery} 的约定给出
\[
\frac{T_{1,(1)}(m)}{f^{(1)}}=T_{1,(1)}(m)=h_1(D_m)=s_{(1)}(D_m).
\]
因此对主核求和中出现的全部分拆都成立
\[
\frac{T_{|\lambda|,\lambda}(m)}{f^\lambda}=s_\lambda(D_m).
\]
故
\[
\mathcal K_m(\mathbf y)
=
\sum_{\lambda}s_\lambda(D_m)s_\lambda(\mathbf y),
\]
其中求和遍历全部分拆并把空分拆的 \(1\) 项包含在内。由经典 Cauchy 恒等式，
\[
\sum_{\lambda}s_\lambda(D_m)s_\lambda(\mathbf y)
=
\prod_{x\in X_m}\prod_{j\ge 1}(1-d_m(x)y_j)^{-1},
\]
得到第一条等式。

另一方面，由命题 \ref{prop:pom-schur-trace-jacobi-trudi-euler-primitive-fiber}，
\[
H_m(z)=\prod_{x\in X_m}(1-d_m(x)z)^{-1}.
\]
对每个 \(j\) 把 \(z\) 替换为 \(y_j\)，再对所有 \(j\ge 1\) 相乘，即得
\[
\prod_{x\in X_m}\prod_{j\ge 1}(1-d_m(x)y_j)^{-1}
=
\prod_{j\ge 1}H_m(y_j).
\]
这就证明了主核恒等式，从而第 \(1\) 条成立。

对第 \(2\) 条，再由命题 \ref{prop:pom-schur-trace-jacobi-trudi-euler-primitive-fiber} 的 Jacobi--Trudi 公式，
\[
s_\lambda(D_m)=\det\bigl(h_{\lambda_i-i+j}(D_m)\bigr)_{1\le i,j\le \ell(\lambda)}.
\]
把
\[
h_r(D_m)=T_{r,(r)}(m)
\qquad
(r\ge 1)
\]
代入，并约定 \(h_0=1\)、\(h_r=0\)（\(r<0\)），便得到一个仅依赖于 \(\lambda\) 的整系数多项式 \(\mathfrak S_\lambda\) 使
\[
\frac{T_{q,\lambda}(m)}{f^\lambda}=s_\lambda(D_m)
=
\mathfrak S_\lambda\bigl(T_{1,(1)}(m),\dots,T_{q,(q)}(m)\bigr).
\]

最后，把 \(\mathbf y\) 特化为 \((z,0,0,\dots)\)。此时经典 Schur 特化给出
\[
s_\lambda(z,0,0,\dots)=0
\quad\text{除非}\quad
\lambda=(q),
\]
而对单行分拆有
\[
s_{(q)}(z,0,0,\dots)=z^q.
\]
因此
\[
\mathcal K_m(z,0,0,\dots)
=
1+\sum_{q\ge 1}\frac{T_{q,(q)}(m)}{f^{(q)}}z^q
=
1+\sum_{q\ge 1}T_{q,(q)}(m)z^q
=
H_m(z),
\]
因为 \(f^{(q)}=1\)。证毕。
\end{proof}

\begin{theorem}[Schur 通道签名对分拆 Gibbs 律的有限层完全可逆性]\label{thm:conclusion-schur-gibbs-signature-complete-invertibility}
\leanverified{paper\_conclusion\_schur\_gibbs\_signature\_complete\_invertibility}
固定整数 \(q\ge 2\) 与分辨率 \(m\)。记
\[
p_{m,q}(\nu):=P^{\mathrm{part}}_{m,q}(\nu)
\qquad(\nu\vdash q),
\]
以及
\[
s_{m,q}(\lambda):=
\frac{T_{q,\lambda}(m)}{f^\lambda\,T_{q,(q)}(m)}
\qquad(\lambda\vdash q).
\]
则映射
\[
p_{m,q}=(p_{m,q}(\nu))_{\nu\vdash q}
\longmapsto
s_{m,q}=(s_{m,q}(\lambda))_{\lambda\vdash q}
\]
是线性同构。更精确地，
\[
s_{m,q}(\lambda)=\sum_{\nu\vdash q}p_{m,q}(\nu)\chi^\lambda(\nu),
\]
并且存在显式反演公式
\[
p_{m,q}(\mu)=\frac1{z_\mu}\sum_{\lambda\vdash q}\chi^\lambda(\mu)\,s_{m,q}(\lambda)
\qquad(\mu\vdash q).
\]
\end{theorem}
\begin{proof}
由正文中分拆 Gibbs 律与 Schur 通道序参量的定义，
\[
m_\lambda(m):=\mathbb E_{m,q}\chi^\lambda(\sigma)
=
\sum_{\nu\vdash q}p_{m,q}(\nu)\chi^\lambda(\nu).
\]
另一方面，
\[
\frac{T_{q,\lambda}(m)}{T_{q,(q)}(m)}=f^\lambda m_\lambda(m),
\]
故
\[
s_{m,q}(\lambda)=m_\lambda(m)
=
\sum_{\nu\vdash q}p_{m,q}(\nu)\chi^\lambda(\nu).
\]

再用对称群字符的列正交关系
\[
\sum_{\lambda\vdash q}\chi^\lambda(\mu)\chi^\lambda(\nu)=z_\mu\,\mathbf 1_{\mu=\nu},
\]
便有
\[
\sum_{\lambda\vdash q}\chi^\lambda(\mu)s_{m,q}(\lambda)
=
\sum_{\nu\vdash q}p_{m,q}(\nu)
\sum_{\lambda\vdash q}\chi^\lambda(\mu)\chi^\lambda(\nu)
=
z_\mu\,p_{m,q}(\mu),
\]
即得所示反演公式。由于正反公式都显式存在，该变换是线性同构。证毕。
\end{proof}

\begin{theorem}[字符签名的 Parseval 等距与相图单纯形]\label{thm:conclusion-schur-signature-parseval-simplex}
在分拆函数空间上定义加权 Hilbert 结构
\[
\langle a,b\rangle_z:=\sum_{\nu\vdash q}z_\nu\,a(\nu)\overline{b(\nu)},
\]
并定义变换
\[
\mathcal F_q:a\longmapsto \widehat a,
\qquad
\widehat a(\lambda):=\sum_{\nu\vdash q}a(\nu)\chi^\lambda(\nu).
\]
则有严格等距律
\[
\sum_{\lambda\vdash q}\widehat a(\lambda)\overline{\widehat b(\lambda)}
=
\sum_{\nu\vdash q}z_\nu\,a(\nu)\overline{b(\nu)}.
\]
特别地，对任意两个分拆概率向量 \(p,p'\) 及其 Schur 签名 \(s,s'\)，有
\[
\sum_{\lambda\vdash q}|s(\lambda)-s'(\lambda)|^2
=
\sum_{\nu\vdash q}z_\nu\,|p(\nu)-p'(\nu)|^2.
\]

再记
\[
v_\nu:=(\chi^\lambda(\nu))_{\lambda\vdash q}\in \CC^{\mathcal P(q)}.
\]
则所有可能的签名恰构成字符列向量的仿射单纯形
\[
\Sigma_q:=\operatorname{conv}\{v_\nu:\nu\vdash q\},
\]
并且任意签名
\[
s=\sum_{\nu\vdash q}p(\nu)v_\nu
\]
中的 \(p(\nu)\) 正是点 \(s\) 在单纯形 \(\Sigma_q\) 中的重心坐标。
\end{theorem}
\begin{proof}
由列正交关系，
\[
\sum_{\lambda\vdash q}\widehat a(\lambda)\overline{\widehat b(\lambda)}
=
\sum_{\mu,\nu\vdash q}
a(\mu)\overline{b(\nu)}
\sum_{\lambda\vdash q}\chi^\lambda(\mu)\chi^\lambda(\nu)
=
\sum_{\nu\vdash q}z_\nu\,a(\nu)\overline{b(\nu)}.
\]
取 \(a=b\) 即得等距公式；再令 \(a=p-p'\)，便得对概率向量的 Parseval 恒等式。

另一方面，由定理 \ref{thm:conclusion-schur-gibbs-signature-complete-invertibility}，
\[
s(\lambda)=\sum_{\nu\vdash q}p(\nu)\chi^\lambda(\nu),
\]
亦即
\[
s=\sum_{\nu\vdash q}p(\nu)v_\nu.
\]
由于 \(\mathcal F_q\) 可逆，列向量族 \(\{v_\nu\}_{\nu\vdash q}\) 线性无关，从而仿射无关；因此其凸包是一个仿射单纯形。重心坐标的唯一性正是上述可逆性的几何表达。证毕。
\end{proof}

\begin{theorem}[压力扇形在字符单纯形中的整胞塌缩]\label{thm:conclusion-pressure-fan-collapses-to-character-simplex-cells}
对每个压力向量 \(p=(P(1),\dots,P(q))\)，定义极大分拆集
\[
R_{\max}(p):=\{\nu\vdash q:\ P(\nu)=P_{\max}\}.
\]
并令
\[
\mathfrak S_q(p):=\lim_{m\to\infty}s_{m,q},
\]
其中极限按正文中的 pressure fan 定理存在。对任意非空集合 \(R\subset \mathcal P(q)\)，若胞腔
\[
\mathcal C_R:=\{p:\ R_{\max}(p)=R\}
\]
非空，则 \(\mathfrak S_q\) 在整块 \(\mathcal C_R\) 上为常值，并且
\[
\mathfrak S_q(p)=\sum_{\nu\in R}\omega_R(\nu)\,v_\nu,
\qquad
\omega_R(\nu):=
\frac{K(\nu)/z_\nu}{\sum_{\mu\in R}K(\mu)/z_\mu}.
\]
特别地：
\begin{enumerate}
\item 若 \(R=\{\nu^\ast\}\) 为单点，则
\[
\mathfrak S_q(p)=v_{\nu^\ast};
\]
\item 若 \(|R|\ge 2\)，则 \(\mathfrak S_q(p)\) 落在面
\[
\operatorname{conv}\{v_\nu:\nu\in R\}
\]
的相对内部。
\end{enumerate}
换言之，压力向量在 Schur 极限层面只留下“哪一组分拆达到最大值”这一组合标签，而不再留下胞腔内部的连续参数。
\end{theorem}
\begin{proof}
若 \(R_{\max}(p)=\{\nu^\ast\}\)，则正文中的纯冻结定理给出
\[
p_{m,q}\to \delta_{\nu^\ast}
\]
指数收敛。代入定理 \ref{thm:conclusion-schur-gibbs-signature-complete-invertibility} 的有限层表示，得
\[
s_{m,q}\to v_{\nu^\ast}.
\]

若 \(|R_{\max}(p)|\ge 2\)，则正文中的多极大分拆定理给出
\[
p_{m,q}(\nu)\to
\begin{cases}
\omega_R(\nu),& \nu\in R,\\
0,& \nu\notin R,
\end{cases}
\]
其中
\[
\omega_R(\nu)=\frac{K(\nu)/z_\nu}{\sum_{\mu\in R}K(\mu)/z_\mu}.
\]
再次代入定理 \ref{thm:conclusion-schur-gibbs-signature-complete-invertibility}，便得
\[
\mathfrak S_q(p)=\sum_{\nu\in R}\omega_R(\nu)v_\nu.
\]
由于每个 \(\omega_R(\nu)\) 都严格为正且总和为 \(1\)，故该点位于相应面之相对内部。右端只由 \(R\) 决定，因此 \(\mathfrak S_q\) 在整块 \(\mathcal C_R\) 上为常值。证毕。
\end{proof}

\begin{corollary}[唯一冻结相的指数可判定性]\label{cor:conclusion-pure-phase-nearest-vertex-exponential-decodability}
若 \(R_{\max}(p)=\{\nu^\ast\}\)，并记
\[
\Delta_q:=\min_{\mu\neq \nu}|v_\mu-v_\nu|_2>0,
\]
则存在常数 \(C_q,\delta_q>0\)，使得
\[
|s_{m,q}-v_{\nu^\ast}|_2\le C_q e^{-m\delta_q}.
\]
从而当 \(m\) 充分大时，\(\nu^\ast\) 是与 \(s_{m,q}\) 距离最近的唯一顶点。亦即，仅凭有限层 Schur 签名的最近顶点判别，即可最终精确恢复冻结主循环型。
\end{corollary}
\begin{proof}
由定理 \ref{thm:conclusion-pressure-fan-collapses-to-character-simplex-cells} 的纯相情形与正文中的指数收敛率，可得每个坐标都有
\[
s_{m,q}(\lambda)=\chi^\lambda(\nu^\ast)+O(e^{-m\delta^\ast}).
\]
由于分拆集合有限，故存在 \(C_q,\delta_q>0\) 使
\[
|s_{m,q}-v_{\nu^\ast}|_2\le C_q e^{-m\delta_q}.
\]
又 \(\{v_\nu\}_{\nu\vdash q}\) 为有限集，故最小点间距 \(\Delta_q\) 严正。若
\[
|s_{m,q}-v_{\nu^\ast}|_2<\frac{\Delta_q}{2},
\]
则三角不等式保证 \(v_{\nu^\ast}\) 为唯一最近顶点。证毕。
\end{proof}

\begin{theorem}[hook 通道生成函数的闭式因子化]\label{thm:conclusion-hook-channel-generating-polynomial-factorization}
对分拆
\[
\nu=1^{c_1}2^{c_2}\cdots q^{c_q}\vdash q,
\]
定义 hook 角色生成多项式
\[
H_\nu(t):=\sum_{k=0}^{q-1}\chi^{(q-k,1^k)}(\nu)\,t^k.
\]
则有严格恒等式
\[
H_\nu(t)=\frac{1}{1+t}\prod_{r=1}^{q}\bigl(1-(-t)^r\bigr)^{c_r}.
\]
并且映射
\[
\nu\longmapsto H_\nu(t)
\]
是单射。更具体地，若写
\[
(1+t)H_\nu(t)=\prod_{d=1}^{q}\Phi_d(-t)^{\,b_d},
\]
其中 \(\Phi_d\) 为第 \(d\) 个 cyclotomic 多项式，则
\[
b_d=\sum_{r:\ d\mid r}c_r,
\qquad
c_r=\sum_{k\ge 1}\mu(k)\,b_{kr}.
\]
\end{theorem}
\begin{proof}
令 \(V=\CC^q\) 为 \(S_q\) 的置换表示，\(W=V/\CC\mathbf 1\) 为标准表示。经典外幂恒等式给出
\[
\sum_{k=0}^{q-1}\chi_{\wedge^k W}(\nu)\,t^k
=
\det(I+t\,\sigma|_W),
\]
其中 \(\sigma\) 为循环型 \(\nu\) 的代表元。又
\[
\wedge^k W\cong S^{(q-k,1^k)},
\]
故左端正是 \(H_\nu(t)\)。

再由 \(V\cong \CC\mathbf 1\oplus W\)，有
\[
\det(I+t\,\sigma|_V)=(1+t)\det(I+t\,\sigma|_W).
\]
而在置换表示上，一个 \(r\)-循环贡献的特征多项式为
\[
\det(I+t\,\sigma_r)=1-(-t)^r.
\]
于是
\[
\det(I+t\,\sigma|_V)=\prod_{r=1}^{q}\bigl(1-(-t)^r\bigr)^{c_r},
\]
两式相除便得
\[
H_\nu(t)=\frac{1}{1+t}\prod_{r=1}^{q}\bigl(1-(-t)^r\bigr)^{c_r}.
\]

再把右端作 cyclotomic 分解：
\[
1-(-t)^r=\prod_{d\mid r}\Phi_d(-t).
\]
故
\[
(1+t)H_\nu(t)=\prod_{d=1}^{q}\Phi_d(-t)^{\,b_d},
\qquad
b_d=\sum_{r:\ d\mid r}c_r.
\]
由 Möbius 反演立得
\[
c_r=\sum_{k\ge 1}\mu(k)\,b_{kr}.
\]
因此 \(H_\nu(t)\) 唯一决定全部 \(c_r\)，从而唯一决定分拆 \(\nu\)。证毕。
\end{proof}

\begin{theorem}[纯相极限中非平凡 hook 通道的无损恢复性]\label{thm:conclusion-pure-phase-hook-family-complete-recovery}
定义 hook 归一化签名
\[
h_{m,q}(k):=
\frac{T_{q,(q-k,1^k)}(m)}
{f^{(q-k,1^k)}\,T_{q,(q)}(m)}
\qquad(0\le k\le q-1),
\]
以及其生成多项式
\[
H_{m,q}(t):=\sum_{k=0}^{q-1}h_{m,q}(k)\,t^k.
\]
则有精确有限层恒等式
\[
H_{m,q}(t)=\sum_{\nu\vdash q}p_{m,q}(\nu)\,H_\nu(t).
\]

若进一步 \(R_{\max}(p)=\{\nu^\ast\}\)，则存在 \(\delta^\ast>0\) 使得
\[
H_{m,q}(t)=H_{\nu^\ast}(t)+O(e^{-m\delta^\ast})
\]
按系数逐项指数收敛。因而当 \(m\) 充分大时，仅凭
\[
\bigl(h_{m,q}(1),\dots,h_{m,q}(q-1)\bigr)
\]
这 \(q-1\) 个非平凡 hook 通道，便可唯一恢复 \(\nu^\ast\)。
\end{theorem}
\begin{proof}
由定理 \ref{thm:conclusion-schur-gibbs-signature-complete-invertibility}，
\[
h_{m,q}(k)=\sum_{\nu\vdash q}p_{m,q}(\nu)\chi^{(q-k,1^k)}(\nu).
\]
对 \(k\) 求和即得
\[
H_{m,q}(t)=\sum_{\nu\vdash q}p_{m,q}(\nu)H_\nu(t).
\]

若 \(R_{\max}(p)=\{\nu^\ast\}\)，则正文中的纯冻结结果给出
\[
p_{m,q}\to \delta_{\nu^\ast}
\]
指数收敛，故
\[
H_{m,q}(t)\to H_{\nu^\ast}(t)
\]
亦按系数指数收敛。

再由定理 \ref{thm:conclusion-hook-channel-generating-polynomial-factorization}，映射 \(\nu\mapsto H_\nu(t)\) 为单射；而分拆集合有限，故不同模板间存在正最小距离。于是当 \(m\) 充分大时，\(H_{m,q}(t)\) 的最近模板必唯一等于 \(H_{\nu^\ast}(t)\)。又因为
\[
h_{m,q}(0)=\frac{T_{q,(q)}(m)}{T_{q,(q)}(m)}=1
\]
恒等成立，故生成多项式已由 \(q-1\) 个非平凡系数 \((h_{m,q}(1),\dots,h_{m,q}(q-1))\) 唯一决定，从而 \(\nu^\ast\) 也被唯一恢复。证毕。
\end{proof}

\begin{conjecture}[hook-stratum 分离]\label{conj:conclusion-hook-stratum-separation}
对每个固定 \(q\)，设
\[
\mathcal R_q:=\{R_{\max}(p):\ p\in \RR^q\}
\]
为全部可实现的极大分拆集。对每个 \(R\in\mathcal R_q\)，定义其 hook 极限模板
\[
\mathcal H_R(t):=
\sum_{\nu\in R}
\frac{K(\nu)/z_\nu}{\sum_{\mu\in R}K(\mu)/z_\mu}\,
H_\nu(t).
\]
则应有
\[
R\neq R'
\quad\Longrightarrow\quad
\mathcal H_R(t)\neq \mathcal H_{R'}(t).
\]
等价地，尽管 hook 投影在整个分拆单纯形上并不保持全局单射，但在真实出现的压力扇形胞腔上，它仍应保持层标签的完全可分。
\end{conjecture}

\noindent
以上结果表明，Schur 通道并非冻结相的附属观测，而是分拆 Gibbs 律的精确 Fourier 坐标：有限层上，分拆概率向量与 Schur 签名严格可逆；极限层上，整个 pressure fan 在字符单纯形中整胞塌缩为顶点与面心；而在纯相极限中，整套不可约谱又进一步坍缩到非平凡 hook 子族而仍保持主循环型的无损可识别性。于是，分拆 Gibbs 几何、Schur 断层坐标与 hook 最小通道族三者在结论层被压接为同一条刚性链。

\endinput


\begin{theorem}[严格分拆间隙下的 Schur 方差非钩化]\label{thm:conclusion-strict-gap-forces-nonhook-schur-variance}
\leanverified{paper\_conclusion\_strict\_gap\_forces\_nonhook\_schur\_variance}
固定整数 \(q\ge 2\)，并沿用正文中循环型单项式 \(\mathcal C_\sigma(m)\)、平均值
\[
\overline{\mathcal C}(m):=\frac{1}{q!}\sum_{\sigma\in S_q}\mathcal C_\sigma(m)
\]
与 Schur 通道 \(T_{q,\lambda}(m)\) 的记号。在 Perron 纯性假设下，若
\[
\Delta_{\mathrm{part}}(q)>0,
\]
则对每个 \(m\ge 1\) 都有精确缩减
\[
\frac{1}{q!}\sum_{\sigma\in S_q}\bigl|\mathcal C_\sigma(m)-\overline{\mathcal C}(m)\bigr|^2
=
\sum_{\substack{\lambda\vdash q,\ \lambda\neq (q)\\ \lambda\ \text{非钩型}}}
\left|\frac{T_{q,\lambda}(m)}{f^\lambda}\right|^2.
\]
其中“非钩型”意指
\[
\lambda\neq (q-k,1^k)
\qquad
\text{对一切 }1\le k\le q-1.
\]
因此，在严格分拆间隙区中，所有非平凡钩型通道在二阶能量上都被完全剥离；剩余非均匀性只能寄居于真正高阶的非钩分拆。
\end{theorem}
\begin{proof}
由定理 \ref{thm:pom-schur-variance-decomposition}，
\[
\frac{1}{q!}\sum_{\sigma\in S_q}\bigl|\mathcal C_\sigma(m)-\overline{\mathcal C}(m)\bigr|^2
=
\sum_{\lambda\vdash q,\ \lambda\neq(q)}
\left|\frac{T_{q,\lambda}(m)}{f^\lambda}\right|^2.
\]
另一方面，定理 \ref{thm:pom-partition-gap-positive-kills-hooks} 断言：若 \(\Delta_{\mathrm{part}}(q)>0\)，则对每个非平凡钩型分拆
\[
\lambda=(q-k,1^k),
\qquad
1\le k\le q-1,
\]
均有
\[
T_{q,\lambda}(m)=0
\qquad
(m\ge 1).
\]
因此方差分解中的全部非平凡钩型项都可精确删去，只剩所示非钩型求和。证毕。
\end{proof}

\begin{corollary}[非均匀性的最小承载复杂度跃迁]\label{cor:conclusion-nonuniformity-minimal-carrier-jumps-to-nonhook}
在定理 \ref{thm:conclusion-strict-gap-forces-nonhook-schur-variance} 的设定下，若对某个 \(m\) 有
\[
\sigma\longmapsto \mathcal C_\sigma(m)
\quad\text{非常值},
\]
则必存在某个分拆 \(\lambda\vdash q\) 满足
\[
\lambda\neq(q),
\qquad
\lambda\ \text{非钩型},
\qquad
T_{q,\lambda}(m)\neq 0.
\]
等价地，任一在严格分拆间隙区仍然存活的非均匀性，其最小承载分拆必满足 Durfee 秩至少为 \(2\)。
\end{corollary}
\begin{proof}
由推论 \ref{cor:pom-schur-nontrivial-channel-activation}，
\[
\sigma\longmapsto \mathcal C_\sigma(m)\ \text{非常值}
\iff
\exists\,\lambda\vdash q,\ \lambda\neq(q),\ T_{q,\lambda}(m)\neq 0.
\]
再由定理 \ref{thm:conclusion-strict-gap-forces-nonhook-schur-variance}，所有非平凡钩型通道都已恒为零，故该 \(\lambda\) 只能是非钩型。Young 图论中，分拆为钩型当且仅当其 Durfee 秩为 \(1\)；故非钩型恰等价于 Durfee 秩至少为 \(2\)。证毕。
\end{proof}

\begin{theorem}[钩型存活所强制的顶层共振多重性]\label{thm:conclusion-hook-survival-forces-top-resonance-multiplicity}
在定理 \ref{thm:pom-hook-channel-implies-partition-gap-closed} 的 Perron 纯性设定下，记顶层共振分拆族
\[
\mathcal R_q:=
\left\{
\nu\vdash q:\ \sum_{r\ge 1}c_r(\nu)P(r)=P(q)
\right\}.
\]
若存在非平凡钩型
\[
\lambda=(q-k,1^k),
\qquad
1\le k\le q-1,
\]
满足
\[
B_{q,\lambda}\neq 0,
\]
则必有：
\begin{enumerate}
\normalfont
  \item \(\Delta_{\mathrm{part}}(q)=0\)；
  \item \(\mathcal R_q\setminus\{(q)\}\neq\varnothing\)；
  \item 存在严格代数消去恒等式
  \[
  \frac{\chi^\lambda((q))}{q}\,\kappa_q
  +
  \sum_{\nu\in\mathcal R_q\setminus\{(q)\}}
  \frac{\chi^\lambda(\nu)}{z_\nu}\prod_{r\ge 1}\kappa_r^{c_r(\nu)}
  =0.
  \]
\end{enumerate}
因此，钩型通道的存活并不只要求“没有严格间隙”；它更强地要求顶层共振不能是 \((q)\) 的单峰独占，而必须出现至少一个非平凡顶层分拆，并与平凡循环项发生精确的代数级抵消。
\end{theorem}
\begin{proof}
第一条正是定理 \ref{thm:pom-hook-channel-implies-partition-gap-closed}。
第三条则正是定理 \ref{thm:pom-hook-channel-leading-term-cancellation} 的结论。

只需证明第二条。若反设
\[
\mathcal R_q=\{(q)\},
\]
则第三条中的求和项为空，从而恒等式退化为
\[
\frac{\chi^\lambda((q))}{q}\,\kappa_q=0.
\]
但对钩型分拆，由 Murnaghan--Nakayama 公式，
\[
\chi^\lambda((q))=(-1)^k\neq 0,
\]
且 \(\kappa_q>0\)。故上式不可能成立，矛盾。于是
\[
\mathcal R_q\setminus\{(q)\}\neq\varnothing.
\]
证毕。
\end{proof}

\begin{conjecture}[torsion 平坦与 strict gap 的双层刚性]\label{conj:conclusion-torsion-flatness-strict-gap-double-rigidity}
固定整数 \(q\ge 2\)。设相关 cyclic 非 Perron 因子在模 \(q\) 上满足 torsion 对称
\[
F_q(t)=G_q(t^q),
\]
并且分拆间隙严格为正：
\[
\Delta_{\mathrm{part}}(q)>0.
\]
则应有严格指数压低
\[
\limsup_{m\to\infty}
\frac{1}{m}
\log\!\left[
\frac{1}{q!}\sum_{\sigma\in S_q}\bigl|\mathcal C_\sigma(m)-\overline{\mathcal C}(m)\bigr|^2
\right]
<
2P(q).
\]
等价地，在这两个条件同时成立时，全部非平凡 Schur 通道都应严格落到 Perron 平方门槛之下；方差主指数不再可能达到 \(2P(q)\)。
\end{conjecture}


\begin{conjecture}[fixed-resolution 的 prime-slice--Schur 最小完备对象应为二元核]\label{conj:conclusion-partition-primeslice-schur-binary-minimal-complete-object}
固定分辨率 \(m\) 与有限集合大小 \(n\)。若某个审计对象要 simultaneously 完备承载
\begin{enumerate}
  \item 分割格上的 meet--\(\gcd\) squarefree 编码热力学；
  \item 纤维多重度多重集 \(\{d_m(x):x\in X_m\}\)；
  \item 全部 Schur 通道 \(\{T_{q,\lambda}(m)\}_{q\ge 1,\ \lambda\vdash q}\)；
  \item 以及压力扇形在 Schur 极限层面的面胞标签；
\end{enumerate}
则其最小完备对象应当不是完整的 prime-slice 编码塔与 Schur 通道塔，而应当恰是二元数据
\[
(\eta,H_m),
\]
其中
\[
\eta:E_n\overset{\sim}{\longrightarrow}S
\]
是定理 \ref{thm:conclusion-partition-primeslice-minimal-uniqueness} 中唯一的边--素数标记，而
\[
H_m(z)=1+\sum_{q\ge 1}T_{q,(q)}(m)z^q
\]
是定理 \ref{thm:conclusion-schur-master-kernel-cauchy-externalization} 的单变量主核。

等价地，在 fixed-resolution 层面，prime-slice 侧与 Schur 侧都应各自存在一个不可再压缩的单生成核，而其余高维对象仅是这两个核的派生阴影。
\end{conjecture}

\noindent
以上结果把三条原本分散的链条压缩为同一组极小性、完备性与相图塌缩原则：
\[
\text{极小 meet--\(\gcd\) 外置}
\Longrightarrow
\text{边--素数重命名外的唯一性}
\Longrightarrow
\text{固定分割的有限温配分函数}
\Longrightarrow
\text{零温唯一基态},
\]
以及
\[
T_{q,(q)}(m)=h_q(D_m)
\Longrightarrow
\{d_m(x)\}_{x\in X_m}\ \text{可完全恢复}
\Longrightarrow
\{T_{q,\lambda}(m)\}\ \text{全部被唯一决定}
\Longrightarrow
\mathcal K_m(\mathbf y)=\prod_{j\ge 1}H_m(y_j).
\]
再者，
\[
p_{m,q}
\Longleftrightarrow
s_{m,q}
\Longrightarrow
\Sigma_q\ \text{中的字符单纯形坐标}
\Longrightarrow
\text{pressure fan 在整胞上塌缩}
\Longrightarrow
\text{纯相中由非平凡 hook 子族无损恢复主循环型}
\Longrightarrow
\text{strict gap 下所有二阶残差被迫转入非钩 Young 几何}.
\]
因此，在分割格侧，最优 prime-slice 编码已被刚性化为唯一模型；在 Schur 主核侧，整个高维断层塔由单变量 complete-symmetric 生成核完全控制；而在 partition Gibbs 几何侧，Schur 通道又给出一组精确 Fourier 坐标，使压力扇形在极限层面直接坍缩为字符单纯形中的顶点与面心。再把 strict partition gap 叠加上去，所有非平凡 hook 通道都会在二阶能量上整体剥离，于是残余非均匀性只能寄居在 Durfee 秩至少为 \(2\) 的非钩分拆中；反过来，任何钩型存活都被迫伴随顶层共振的非单峰多重性与精确代数消去。三条链在结论层因而收束为同一类现象：一旦达到最小坐标预算，剩余结构不再表现为自由选择，而表现为由极小对象唯一生成的全塔闭包、整胞相图，以及经双层剥离后才暴露出来的高阶残差几何。

\endinput
